\section{The \sys{} Multi-Agent Fleet}
\label{sec:methods}

\sys{} is not a single agent. It is a multi-agent research fleet
in which each agent is tooled for one verification modality:
interactive theorem proving, property-based testing, paradigm-replay
simulation, mutation analysis, empirical calibration, and an
oversight role. Frontier large language models drive the agents.
Two architectural properties are essential.

\paragraph{Verification-gated handoffs.}
Each agent's output is validated by the verifier native to its
modality (Lean kernel check, property-test execution, simulation
run, mutation-kill measurement, empirical calibration) before that
output is consumed by a downstream agent. Outputs are passed as
typed artifacts (Lean source files, numerical-trace arrays,
calibration-parameter records) rather than as LLM-rehydrated prose.
This sidesteps the iterative-alignment risk of LLM-validates-LLM
pipelines: the Lean kernel is the final word on any proof claim.

\paragraph{Persistent shared context.}
Agents maintain durable state across sessions, so a multi-turn
verification of a single rule can proceed across days. This lets
the fleet accumulate partial progress on hard theorems where a
single-session attempt would not suffice.

\paragraph{Six verification modalities.}
The fleet covers six modalities, each targeting a distinct failure
mode:
\begin{enumerate}[leftmargin=1.5em,itemsep=0.1em]
\item \textbf{Lean~4 algebraic proof.} Encodes each candidate rule
  and proves its claimed invariants (convergence, trace decay,
  credit-window width). Catches structural errors: wrong convergence
  rate, missing hypothesis, vacuous closure.
\item \textbf{Property-based testing.} Hypothesis (the Python PBT
  library) draws parameter tuples from the modelled regime and shrinks
  any falsifying sample. Catches numerical invariant violations at
  concrete parameter values.
\item \textbf{Paradigm-replay simulation.} Re-implements each rule
  in Brian~2 under the \citet{tang2024dynamic} closed-loop paradigm.
  Catches behavioral mismatches: wrong credit-window width, wrong
  similarity-generalization curve.
\item \textbf{Mutation analysis.} Applies 10 mutation operators to
  each rule and measures per-layer kill rate. Quantifies the
  complementarity of the three verification layers.
\item \textbf{Empirical calibration.} Fits the eligibility-trace
  width to Tang et al.\ Fig.~3c ($n=55$ actions, $4.4\%$ error).
\item \textbf{Oversight.} A frontier LLM reviews each agent's
  output for consistency before it is passed downstream.
\end{enumerate}

\paragraph{Self-evolving loop.}
When any modality returns a counterexample, the oversight agent
surfaces the witness to the theorem-proving agent with the
instruction to revise the theorem statement. The revised statement
is re-submitted to all six modalities. The loop terminates when
all modalities return \texttt{closed} or when a maximum iteration
count is reached. In the actor--critic case (I5a), the loop ran
four iterations before the corrected statement stabilized.

\paragraph{Candidate rules.}
We apply the fleet to four candidate dopamine credit-assignment
rules: TD($\lambda$) with eligibility traces, action-prediction
error (APE), Rescorla--Wagner (following \citet{markowitz2023spontaneous}),
and time-magnitude rank-4 distributional RL (TMRL,
\citet{sousa2025tmrl}). Each rule is encoded in Lean~4 with a
shared substrate: state and action are finite types, rewards are
real-valued, and a policy is a mapping $s \mapsto \mathrm{PMF}(a)$.
